\chapter{Matriz de Rastreabilidade de Requisitos}
\label{cap_matriz_rastreabilidade}

A rastreabilidade de requisitos é a capacidade de acompanhar a vida de um requisito em direção prospectiva e retrospectiva ao longo de todo o ciclo de vida do software \cite{sommerville2019,iso29148}. Em conformidade com a norma ISO/IEC/IEEE 29148:2018, a rastreabilidade assegura que:
\begin{enumerate}
    \item Cada necessidade clínica e pedagógica esteja concretizada em ao menos um Requisito Funcional (RF);
    \item Cada RF esteja mapeado a casos de uso do sistema, atores responsáveis e entidades do modelo de dados;
    \item As restrições de qualidade expressas nos Requisitos Não Funcionais (RNFs) incidam com exatidão sobre os módulos funcionais correspondentes;
    \item Modificações futuras em regras de negócio possam ser analisadas quanto ao seu impacto arquitetural antes de serem implementadas.
\end{enumerate}

% ===========================================================================
\section{Mapeamento Bidirecional de Requisitos}
\label{sec_mapeamento_bidirecional}
% ===========================================================================

As Tabelas \ref{tab_matriz_rastreabilidade_1} e \ref{tab_matriz_rastreabilidade_2} consolidam a Matriz de Rastreabilidade Bidirecional do SIGEA-GTT, correlacionando os 25 Requisitos Funcionais com seus respectivos atores primários, casos de uso operacionais, requisitos não funcionais incidentes e tabelas relacionais do banco de dados PostgreSQL.

\begin{table}[htb]
\centering
\small
\caption{Matriz de Rastreabilidade Bidirecional (RF01 a RF13).}
\label{tab_matriz_rastreabilidade_1}
\begin{tabularx}{\textwidth}{l c >{\raggedright\arraybackslash}p{3.5cm} >{\raggedright\arraybackslash}X >{\raggedright\arraybackslash}p{3.8cm}}
\toprule
\textbf{RF} & \textbf{Atores} & \textbf{Caso de Uso (UC)} & \textbf{RNF Associados} & \textbf{Entidades / Tabelas} \\
\midrule
\textbf{RF01} & Todos & UC01 (Autenticar) & RNF03, RNF07 & \texttt{users} \\
\midrule
\textbf{RF02} & Todos & UC01 (Autenticar) & RNF03, RNF07 & \texttt{users} \\
\midrule
\textbf{RF03} & Todos & UC02 (Alterar Senha) & RNF03, RNF07 & \texttt{users} \\
\midrule
\textbf{RF04} & ADMIN & UC04 (Gerenciar Usuários) & RNF02, RNF05, RNF08, RNF09 & \texttt{users} \\
\midrule
\textbf{RF05} & ADMIN & UC04 (Gerenciar Usuários) & RNF03, RNF09, RNF10 & \texttt{users} \\
\midrule
\textbf{RF06} & ADMIN & UC05 (CRUD Módulos GTT) & RNF02, RNF05, RNF08 & \texttt{modulos\_gtt} \\
\midrule
\textbf{RF07} & ADMIN & UC06 (CRUD Gatilhos IHI) & RNF02, RNF05, RNF08 & \texttt{gatilhos\_gtt} \\
\midrule
\textbf{RF08} & Todos & UC03 (Consultar Catálogo) & RNF01, RNF06 & \texttt{modulos\_gtt}, \texttt{gatilhos\_gtt} \\
\midrule
\textbf{RF09} & Todos & UC03 (Guia MERP) & RNF01, RNF06 & \texttt{gravidades\_merp} \\
\midrule
\textbf{RF10} & ADMIN & UC07 (Gerenciar Turmas) & RNF05, RNF08, RNF09 & \texttt{turmas} \\
\midrule
\textbf{RF11} & ADMIN & UC07 (Gerenciar Turmas) & RNF08, RNF09, RNF10 & \texttt{turmas}, \texttt{users} \\
\midrule
\textbf{RF12} & ADMIN & UC07 (Gerenciar Turmas) & RNF02, RNF05, RNF08 & \texttt{turmas\_alunos}, \texttt{users} \\
\midrule
\textbf{RF13} & ADMIN & UC07 (Gerenciar Turmas) & RNF08, RNF09, RNF10 & \texttt{turmas}, \texttt{submissoes} \\
\bottomrule
\end{tabularx}
\fonte{Elaborado pelo autor (2026).}
\end{table}

\begin{table}[htb]
\centering
\small
\caption{Matriz de Rastreabilidade Bidirecional (RF14 a RF25).}
\label{tab_matriz_rastreabilidade_2}
\begin{tabularx}{\textwidth}{l c >{\raggedright\arraybackslash}p{3.5cm} >{\raggedright\arraybackslash}X >{\raggedright\arraybackslash}p{3.8cm}}
\toprule
\textbf{RF} & \textbf{Atores} & \textbf{Caso de Uso (UC)} & \textbf{RNF Associados} & \textbf{Entidades / Tabelas} \\
\midrule
\textbf{RF14} & PROF & UC09 (Criar Atividade) & RNF05, RNF08, RNF09 & \texttt{atividades} \\
\midrule
\textbf{RF15} & PROF & UC09 (Criar Prontuário) & RNF01, RNF04, RNF06 & \texttt{atividades} (JSONB) \\
\midrule
\textbf{RF16} & STUDENT & UC12 (Cronômetro 20 min) & RNF01, RNF06, RNF07 & \texttt{submissoes} \\
\midrule
\textbf{RF17} & STUDENT & UC13 (Identificar Gatilhos) & RNF01, RNF06, RNF09 & \texttt{submissoes\_gatilhos} \\
\midrule
\textbf{RF18} & STUDENT & UC14 (Ishikawa 6M) & RNF01, RNF06, RNF09 & \texttt{submissoes} (JSONB) \\
\midrule
\textbf{RF19} & STUDENT & UC14 (Matriz GUT) & RNF01, RNF06, RNF09 & \texttt{submissoes} (JSONB) \\
\midrule
\textbf{RF20} & STUDENT & UC14 (Matriz 5W2H) & RNF01, RNF06, RNF09 & \texttt{submissoes} (JSONB) \\
\midrule
\textbf{RF21} & STUDENT & UC14 (Ciclo PDCA) & RNF01, RNF06, RNF09 & \texttt{submissoes} (JSONB) \\
\midrule
\textbf{RF22} & STUDENT & UC14 (SWOT / Brainstorm) & RNF01, RNF06, RNF09 & \texttt{submissoes} (JSONB) \\
\midrule
\textbf{RF23} & STUDENT & UC15 (Submeter Resolução) & RNF02, RNF07, RNF08 & \texttt{submissoes} \\
\midrule
\textbf{RF24} & PROF & UC10 (Avaliar e Emitir Parecer) & RNF01, RNF07, RNF09 & \texttt{submissoes} \\
\midrule
\textbf{RF25} & PROF, ADM & UC08 / UC11 (Dashboards) & RNF01, RNF02, RNF06 & \texttt{submissoes}, \texttt{turmas} \\
\bottomrule
\end{tabularx}
\fonte{Elaborado pelo autor (2026).}
\end{table}

% ===========================================================================
\section{Análise de Cobertura e Impacto}
\label{sec_analise_cobertura}
% ===========================================================================

A inspeção detalhada da matriz de rastreabilidade evidencia que:
\begin{enumerate}
    \item \textbf{Cobertura Funcional Total (100\%)}: Todos os 25 requisitos funcionais possuem ao menos um caso de uso executor correspondente e mapeamento direto para tabelas estruturadas ou campos JSONB de alta performance no PostgreSQL 16;
    \item \textbf{Cobertura de Qualidade Total (100\%)}: Todos os 10 requisitos não funcionais incidem transversalmente sobre os módulos da aplicação, com destaque para a usabilidade e ergonomia hospitalar (\texttt{RNF01} e \texttt{RNF06}), a segurança criptográfica (\texttt{RNF03}) e a confiabilidade atestada pelo Quality Gate (\texttt{RNF09});
    \item \textbf{Isolamento de Domínio e LGPD}: Nenhum requisito funcional manipula bases externas ou dados identificáveis de pacientes reais, cumprindo estritamente a exigência de segregação estabelecida pelo \texttt{RNF04}.
\end{enumerate}
